% Guard: macros.tex is \input both by main.tex (after the packages it needs,
% e.g. amsmath for \DeclareMathOperator) and by individual section files so
% they remain compilable on their own. Loading it twice would redefine every
% \newcommand and abort, so the second and later loads are no-ops.
\ifdefined\PQRmacrosloaded\endinput\fi
\let\PQRmacrosloaded\relax

% --- 0. revision marking ---
% Text added to Section 2 during the current revision is wrapped in \added{...}
% so the advisors can see at a glance what is new. Redefine to \relax to print
% the paper without the marking.
\newcommand{\added}[1]{\textcolor{blue}{#1}}
\long\def\added#1{\begingroup\color{blue}#1\endgroup}  % \long: may span paragraphs


% --- A. operator-style score and function names ---
\newcommand{\HD}{\mathit{HD}}                        % histogram distance; use \HD_\gamma
\newcommand{\Attr}{\mathit{Attr}}                    % attribute set of a schema
\newcommand{\Attrs}{\mathit{Attrs}}                  % attribute set argument
% Abbreviated in MATH only; the prose still spells the names out in full.
% Spelling them out inside formulas makes displayed equations far too wide
% for a two-column layout (the product-rule step below overflowed by 200pt).
\newcommand{\Use}{\mathit{Use}}
\newcommand{\Cov}{\mathit{Cov}}
\newcommand{\Rep}{\mathit{Rep}}
\newcommand{\DivS}{\mathit{Div}}
\newcommand{\TVD}{\mathit{TVD}}
\newcommand{\GS}{\mathit{GS}}                        % global sensitivity
\DeclareMathOperator{\Dom}{Dom}
\DeclareMathOperator{\Lap}{Laplace}

% --- B. core objects ---
\newcommand{\Qone}{Q_1}
\newcommand{\Qtwo}{Q_2}
\newcommand{\pD}{p(D)}
\newcommand{\views}{\mathcal{P}}
\newcommand{\hist}[2]{H_{#1,#2}}                     % \hist{\Qone(\pD)}{A}
\newcommand{\dhist}{\Delta H}
\newcommand{\nhist}{\widetilde{\Delta H}}            % the released noisy object

% --- C. privacy budget and parameters ---
\newcommand{\eps}{\epsilon}
\newcommand{\epsP}{\epsilon_P}
\newcommand{\epsS}{\epsilon_S}
\newcommand{\epsH}{\epsilon_H}
\newcommand{\gammaref}{\gamma_{\mathrm{ref}}}
\newcommand{\gammaeval}{\gamma_{\mathrm{eval}}}
\newcommand{\taumax}{\tau_{\max}}

% --- D. mechanisms ---
\newcommand{\ExpMech}{\textsc{ExpMech}}
\newcommand{\OneShotTop}[1]{\textsc{OneShotTop-}#1}
\newcommand{\AdaptiveTau}{\textsc{AdaptiveTau}}
% --- E. experiment baselines (used in sections/experiments.tex) ---
\newcommand{\Bone}{\ensuremath{\mathsf{B1}}}
\newcommand{\Btwo}{\ensuremath{\mathsf{B2}}}
\newcommand{\Bthree}{\ensuremath{\mathsf{B3}}}

\section{Model and Problem Statement}\label{sec:model}

\subsection{Private Query-Difference Setting}
\label{subsec:private-query-difference-setting}

We consider a trusted server that holds a private database $D$ and receives, or derives, two related queries $Q_1$ and $Q_2$ over the same database.
The goal is to explain the semantic difference between the corresponding query outputs $Q_1(D)$ and $Q_2(D)$ while preserving the privacy of individuals in $D$.

In the query-refinement setting, $Q_1=Q$ is the original user query and $Q_2=Q'$ is a refined query produced by a refinement mechanism so that $Q'(D)$ satisfies a set of output constraints.
In the consistent-query answering setting, the two outputs correspond to evaluating a query directly on the original database and evaluating it under repair semantics.
In both settings, our framework treats the mechanism that produces the related output as a black box and focuses only on explaining the difference between the two resulting outputs.

Concretely, the workflow is the following (Figure~\ref{fig:framework}).
An analyst submits a query $Q_1$ to a trusted server holding the private database $D$; the server obtains the related query $Q_2$, either because the analyst supplies it or because an external mechanism --- a refinement algorithm, a repair procedure --- derives it.
The server then evaluates both queries, privately selects a small set of predicate views and, within them, the explanation candidates that best characterise the difference, and finally releases noisy histograms for the selected candidates.
Only this last object leaves the server; the analyst never observes $D$, the exact query outputs, or the exact scores that drove the selection.
The mechanism producing $Q_2$ is treated as a black box throughout, and its own privacy cost, if any, is outside the scope of the analysis here.
\alon{Done: trusted-server workflow described above and depicted in Figure~\ref{fig:framework}.}

\subsection{Explanation Candidates}
\label{subsec:explanation-candidates}

An explanation in our framework compares the distribution of one attribute across the two query outputs, restricted to the same semantic region of the database.
The semantic region is represented by a predicate view $p(D)$, and the compared distribution is represented by the one-way marginal of an attribute $A$.

\begin{definition}[Histogram Explanation Candidate]
Given a database $D$, two related queries $Q_1,Q_2$, a predicate view $p$, and an attribute $A \in \Attr(D)$, a \emph{histogram explanation candidate} is a pair
\[
    e=(p,A).
\]
The candidate induces the two histograms
\[
    H_{Q_1(p(D)),A}
    \qquad\text{and}\qquad
    H_{Q_2(p(D)),A}.
\]
\end{definition}

The predicate view $p(D)$ localizes the explanation to a semantic region of the data, while the attribute $A$ determines the distribution being compared.
A candidate is informative when the two induced histograms differ substantially, since such a difference indicates that the change from $Q_1$ to $Q_2$ affected the selected population with respect to attribute $A$.

\begin{example}[Explanation candidate]
\label{ex:candidate}
Continuing Example~\ref{ex:running}, take the view $p=(\texttt{education}=\texttt{Masters})$ and the attribute $A=\texttt{occupation}$.
The candidate $e=(p,A)$ induces two histograms over occupations: $\hist{\Qone(\pD)}{A}$, counting Master's-degree holders working at least $50$ hours who are selected by the original query, and $\hist{\Qtwo(\pD)}{A}$, counting those selected by the refined query.
The candidate is informative to the extent that these two occupational profiles differ; if the refinement shifts the population from managerial towards professional-specialty roles, that shift is visible as a difference between the two histograms, and is invisible in the aggregate value alone.
\end{example}
\alon{Done: running example added as Example~\ref{ex:candidate}.}

\subsection{Ideal Explanation Quality}
\label{sec:explanation-quality}

We evaluate explanation candidates according to four desiderata: \emph{interestingness}, \emph{representativeness}, \emph{coverage}, and \emph{diversity}.
Similar families of criteria are commonly used in explanation and data-exploration systems: interestingness captures whether an explanation highlights a meaningful deviation, while diversity encourages a set of explanations to cover complementary aspects of the data rather than repeatedly showing the same phenomenon \cite{FEDEX}.
Histogram-based explanations and distributional summaries are also used in private explanation frameworks such as DPXPlain \cite{TaoGMR22}.
The closest analogue is the recent work on differentially private explanations for clustering results~\cite{DPClustX}, which likewise proceeds by showing that the natural quality measures are too sensitive for direct private use and then constructing low-sensitivity replacements.
Our setting differs in what is being explained --- the difference between two \emph{query outputs} over the same database, rather than the deviation of a cluster from the full dataset --- but the design pressure is the same, and we adopt the same overall strategy.
\alon{Done: cited~\cite{DPClustX}; representativeness and coverage follow the exploration measures of~\cite{FEDEX,ICDE23} already cited above.}

For an attribute $A$, let $\Dom(D[A])$ denote the finite set of values, or buckets, used in the histogram of $A$.
For a candidate $e=(p,A)$, the \emph{interestingness} component measures how much the distribution of $A$ changes between $Q_1(p(D))$ and $Q_2(p(D))$:
\begin{multline*}
I(p,A) = \sum_{v\in\Dom(D[A])} \Big|
\Pr(t[A]=v \mid t \in Q_1(p(D)))\\
-
\Pr(t[A]=v \mid t \in Q_2(p(D)))
\Big|.
\end{multline*}
Intuitively, an explanation is interesting if it highlights an attribute whose distribution is substantially affected by the query change.
When either $Q_1(p(D))$ or $Q_2(p(D))$ is empty the conditional distributions above are undefined; we assign such candidates $I(p,A)=0$ rather than discarding them, so that the domain of the objective does not itself depend on the data.
The same convention applies to the proxy of Definition~\ref{def:hd}, where it arises automatically: an empty query output contributes zero to every bin.
\alon{Done: empty outputs are scored $0$, not discarded, so that the candidate set is data-independent.}

The \emph{representativeness} component favors predicate views that are not overly exceptional or idiosyncratic relative to the underlying database.
We measure how close the distribution of $A$ inside the view $p(D)$ is to its distribution in the full database:
\begin{align}
R(p,A)
&= 1-\TVD\!\left(P_{D,A},\,P_{p(D),A}\right)\\
&= 1-\frac{1}{2}\!\!\sum_{v \in \Dom(D[A])}\!\!
\Big| \Pr(t[A]{=}v \mid t \in D) \nonumber\\[-2pt]
&\hphantom{{}= 1-\frac{1}{2}\!\!\sum_{v \in \Dom(D[A])}\!\!\Big|}
- \Pr(t[A]{=}v \mid t \in p(D)) \Big|.
\label{eq:ideal-representativeness}
\end{align}
This prevents explanations from focusing on small or unusual regions that may exhibit large changes but are not representative of the data.
The factor of $\frac{1}{2}$ is the one carried by $\TVD$ (Definition~\ref{def:tvd}) and is required for $R$ to be well defined as a quality score: without it the sum ranges over $[0,2]$ and $R$ can be as low as $-1$, so a maximally unrepresentative view would be rewarded with a large negative multiplier in $U^*$ below rather than penalised.
\alon{Done: normalisation fixed to $1-\TVD$; the unnormalised form could range below $0$. The low-sensitivity counterpart used by the algorithm is Eq.~\eqref{eq:representativeness-proxy}, which additionally smooths the view-side denominator.}

The \emph{coverage} component favors views that contain sufficient evidence from the compared query outputs.
High coverage makes the explanation statistically more reliable and less likely to be driven by a few tuples:
\[
C(p)
=
|p(D)| \cdot |Q_1(p(D))| \cdot |Q_2(p(D))|.
\]
The product form is deliberate: it vanishes whenever \emph{either} query output is empty inside the view, which is the case in which no comparison between $Q_1$ and $Q_2$ is possible at all, and it grows with the evidence available on both sides.
A sum would not have this property, as it would reward a view that is large on one side and empty on the other.
\alon{Done: kept the product, with the justification that it vanishes when either side is empty; the algorithm's normalised proxy (Eq.~\eqref{eq:coverage-proxy}) is the low-sensitivity counterpart, related by dividing by $|D|$ and dropping the $|p(D)|$ factor.}

For a single candidate, the ideal quality is
\[
U^*(p,A)
=
I(p,A) \cdot R(p,A) \cdot C(p).
\]
When releasing several explanations, the selected set should also be diverse.
Diversity encourages the explanations to expose different semantic aspects of the change from $Q_1(D)$ to $Q_2(D)$, rather than repeatedly reporting the same phenomenon.
We therefore state the ideal quality of a \emph{set} $E$ of candidates, which is what Definition~\ref{def:query-difference-explanation} asks to maximise.
For an attribute $A$, let $U^*_{(1)}(A)\ge U^*_{(2)}(A)\ge\cdots$ be the ideal scores of the candidates of $E$ that use $A$, in decreasing order, and let $w_1\ge w_2\ge\cdots\ge0$ with $w_1=1$ be a non-increasing weight profile.
Then
\begin{equation}
U^*_w(E)=\sum_{A}\;\sum_{j\ge1} w_j\,U^*_{(j)}(A).
\label{eq:ideal-set-utility}
\end{equation}
Within each attribute group the strongest candidate contributes in full and every further candidate is discounted by its rank, so repeatedly explaining the same attribute yields diminishing returns.
We use the graded profile $w_j=\lambda^{\,j-1}$ for $0\le\lambda\le1$, which is the form the algorithm's diversity term takes in Equation~\eqref{eq:diversity-score}; $\lambda=1$ removes the diversity requirement entirely and $\lambda=0$ admits at most one explanation per attribute.
A flat profile $w=(1,\lambda,\lambda,\dots)$, charging every repeat the same $\lambda$, is also admissible and coincides with the graded one on the first two members of a group; the two differ from the third onwards, which matters once $k$ is large enough for an attribute to be used three or more times.
Section~\ref{subsubsec:set-objective} shows that the proxy counterpart of Eq.~\eqref{eq:ideal-set-utility}, obtained by replacing $U^*$ with $\HD_\gamma$, is order-independent and monotone submodular for any such $w$, and Theorem~\ref{thm:non-private-hardness} shows that it is in fact maximised \emph{exactly} by the greedy selection of Algorithm~\ref{alg:topk}.
\alon{Done: set-level utility stated in Eq.~\eqref{eq:ideal-set-utility}. It was previously written as $\lambda\sum_e U^*(e)+(1-\lambda)\sum_A\max U^*$, which expands to the \emph{flat} weights $(1,\lambda,\lambda,\dots)$ and therefore disagreed with the graded profile $\lambda^{j-1}$ used by Eq.~\eqref{eq:diversity-score} and Definition~\ref{def:setobj} from the third member of a group onwards --- so the proxy was not in fact a counterpart of the ideal. Both are now stated in the same general weighted form.}

The ideal objective above captures the explanations we would like to select in the absence of privacy and computational constraints.
However, directly optimizing it is problematic for differentially private selection.
In particular, the score contains conditional distributions and coverage terms that may change significantly when a single tuple is added or removed, especially for small predicate views or small query outputs.
This motivates the low-sensitivity proxy measures introduced in Section~\ref{sec:solution}.

\begin{proposition}[Sensitivity of the ideal quality function]
\label{prop:ideal-utility-sensitivity}
The global sensitivity of $U^*$ is not bounded by any constant.
Specifically, there exist neighbouring databases $D\sim D'$, a view $p$ and an attribute $A$ for which
$|U^*(p,A;D)-U^*(p,A;D')| = \Omega(|D|^2)$.
\end{proposition}

\begin{proof}
We give an explicit family of instances, then note a second, independent source of instability.

\emph{The $\Omega(|D|^2)$ construction.}
Let $D$ consist of $n$ tuples, all satisfying $p$, so that $p(D)=D$ and $|D|=n$.
Give $n-1$ of them $A$-value $a$ and let each satisfy $Q_1$ but not $Q_2$; give the remaining tuple $A$-value $b\ne a$ and let it satisfy $Q_2$ but not $Q_1$.
Then
\[
C(p)=|p(D)|\cdot|Q_1(p(D))|\cdot|Q_2(p(D))| = n(n-1)\cdot 1 ,
\]
the two conditional distributions are point masses on $a$ and on $b$ respectively, so $I(p,A)=2$, its maximum; and $p(D)=D$ gives $\TVD(P_{D,A},P_{p(D),A})=0$, so $R(p,A)=1$.
Hence $U^*(p,A;D)=2n(n-1)$.

Let $D'=D\setminus\{t\}$ where $t$ is the single $Q_2$ tuple.
Now $Q_2(p(D'))=\emptyset$, so $C(p)=0$ and, by the convention fixed above, $I(p,A)=0$; therefore $U^*(p,A;D')=0$.
The two databases are neighbours and
\[
|U^*(p,A;D)-U^*(p,A;D')| = 2n(n-1)=\Omega(|D|^2),
\]
which is the claim.
Note that the three factors are large \emph{simultaneously} here, which is what makes the bound a statement about $U^*$ rather than about $C$ alone.
This is the price of the product form for $C$: exactly the property that makes it vanish when either side is empty makes it enormous when one side is nearly empty.

\emph{A second source: interestingness.}
Even holding $C$ fixed, $I$ is a difference of two \emph{conditional} distributions, so its stability depends on $|Q_1(p(D))|$ and $|Q_2(p(D))|$ rather than on $|D|$.
Let $Q_1(p(D))=\{t\}$ with $t[A]=a$ and let $Q_2(p(D))$ be large with all its mass on a single value $b\neq a$, so $I(p,A)=2$.
Adding to $Q_1(p(D))$ one tuple with $A$-value $b$ moves $I$ to $1$: a change of $1$, independent of $|D|$.
(The added value must lie in the support of $Q_2(p(D))$; adding a value fresh to both outputs leaves $I$ unchanged, so it is the overlap that drives the change.)
Consequently no sensitivity bound for $I$ that decays with $|D|$ is available, and in particular the $1/|D|$-type bounds that the proxies of Section~\ref{sec:solution} enjoy cannot hold for $U^*$.

Both failures share a cause: the quantities are \emph{ratios or products of counts} that can be driven by a handful of tuples.
This is precisely the difficulty that the low-sensitivity proxies of Section~\ref{sec:solution} are designed to remove --- $\Cov$ replaces the product by a sum normalised by the public $|D|$ (Eq.~\eqref{eq:coverage-proxy}), and $\HD_\gamma$ replaces the conditional distributions by counts with a smoothed denominator (Definition~\ref{def:hd}).
\end{proof}
\alon{Done: proof added, showing the $\Omega(|D|^2)$ blow-up of the coverage factor and the $|D|$-independent instability of the interestingness factor.}

\subsection{Problem Definition}
\label{subsec:problem-definition}

We can now state the explanation task using the candidate and quality notions above.
The non-private version asks for high-quality explanations of the difference between two related query outputs.
The private version adds the requirement that the released explanations and histograms satisfy differential privacy with respect to the private database $D$.

\begin{definition}[Query-Difference Explanation Problem]
\label{def:query-difference-explanation}
Given a database $D$, two related queries $Q_1,Q_2$, a set of predicate views $\mathcal{P}$, an attribute set $\Attr(D)$, and an integer $k$, return a set $E$ of $k$ histogram explanation candidates $(p,A)$ that maximizes the ideal explanation quality for the semantic difference between $Q_1(D)$ and $Q_2(D)$.
\end{definition}

\begin{definition}[Differentially Private Query-Difference Explanation Problem]
\label{def:problem}
Given a private database $D$, two related queries $Q_1,Q_2$, a set of predicate views $\mathcal{P}$, an attribute set $\Attr(D)$, a privacy budget $\epsilon>0$, and an integer $k$, output $k$ histogram-based explanations for the difference between $Q_1(D)$ and $Q_2(D)$ while satisfying $\epsilon$-differential privacy with respect to $D$.
\end{definition}

The released object is the noisy \emph{difference} histogram $\nhist_{p,A}$ for each selected candidate.
Releasing the two histograms $\hat H_{Q_1(p(D)),A}$ and $\hat H_{Q_2(p(D)),A}$ separately is also possible, and is strictly more informative, since the difference is recoverable from the pair by post-processing, but it is not free: each histogram must then be noised independently, so under a fixed budget $\epsilon_H$ each carries roughly twice the noise of a single difference release.
We define the problem in terms of the difference histogram and note the paired variant where the distinction matters.
\alon{Done: the released object is fixed to be the noisy difference histogram, with the paired alternative and its cost noted.}

\subsection{Challenges}
\label{subsec:model-challenges}

The ideal problem above is difficult for two reasons.
First, the space of possible predicate views is combinatorial in the number of attributes and values, so enumerating all possible semantic regions is infeasible.
Second, the ideal utility function may have high sensitivity, which would require adding large amounts of noise or would make private selection ineffective.
Therefore, the algorithm in the next section restricts the candidate space to predicate-based views and replaces the ideal utility components with low-sensitivity proxy measures.
